%=============================================================================
% W4i17: EUCLID DR1 PREPARATION --- COMPLETE ANALYSIS BLUEPRINT
% 7-Month Countdown to First Critical DFD Test
% Wave 4, Iteration 17 | 20 March 2026 | Model: Opus 4.6
%=============================================================================

\documentclass[11pt,a4paper]{article}
\usepackage[margin=2.2cm]{geometry}
\usepackage{amsmath,amssymb}
\usepackage{booktabs}
\usepackage{enumitem}
\usepackage{float}
\usepackage{xcolor}
\usepackage[most]{tcolorbox}
\usepackage{hyperref}
\usepackage{longtable}
\usepackage{array}

\newcommand{\LCDM}{$\Lambda$CDM}
\newcommand{\DFD}{DFD}
\newcommand{\ee}[1]{\times 10^{#1}}
\newcommand{\kmu}{k_\mu}
\newcommand{\Geff}{G_{\rm eff}}
\newcommand{\Dpsi}{\Delta\psi}

\newtcolorbox{findbox}[1][]{colback=yellow!5!white, colframe=red!70!black,
  fonttitle=\bfseries, title={#1}}
\newtcolorbox{granitebox}[1][]{colback=green!5,colframe=green!60!black,
  fonttitle=\bfseries,title={#1}}
\newtcolorbox{distbox}[1][]{colback=blue!3,colframe=blue!50!black,
  fonttitle=\bfseries,title={#1}}
\newtcolorbox{contigbox}[1][]{colback=purple!3,colframe=purple!50!black,
  fonttitle=\bfseries,title={#1}}
\newtcolorbox{pipelinebox}[1][]{colback=orange!3,colframe=orange!60!black,
  fonttitle=\bfseries,title={#1}}

\title{\textbf{Wave 4, Iteration 17:\\
Euclid DR1 Preparation --- Complete Analysis Blueprint}\\[0.5em]
\large 7-Month Countdown $\cdot$ Data Product Mapping $\cdot$
DFD Prediction Formulae $\cdot$ Weekend-Viable Pipeline\\[0.3em]
\large Competitive Landscape and Falsification Contingency}

\author{DFD Research Programme --- Agent 16 (Survey Preparation, Opus 4.6)}
\date{20 March 2026}

\begin{document}
\maketitle

%=============================================================================
\section*{Executive Summary: Top 5 Findings}
%=============================================================================

\begin{findbox}[TOP 5 FINDINGS --- WAVE 4, ITERATION 17]

\begin{enumerate}

\item \textbf{CRITICAL --- Euclid DR1 Confirmed 21 October 2026,
  $\sim$1900\,deg$^2$, First Cosmology Release:}
  ESA confirms DR1 on 21 October 2026 covering $\sim$1900\,deg$^2$
  ($\sim$500\,deg$^2$ North, $\sim$1400\,deg$^2$ South). This is the
  first Euclid release containing cosmology-grade data products: VIS
  shear catalog ($\sim$30 galaxies\,arcmin$^{-2}$, $\sigma_e \approx 0.26$),
  NISP photometric redshifts ($0.2 < z < 2.5$), NISP spectroscopic
  redshifts ($0.9 < z < 1.8$, H$\alpha$ emitters, $\sim$1700\,deg$^{-2}$),
  and merged multi-wavelength catalogs. At 1900\,deg$^2$ with 30\,arcmin$^{-2}$,
  this yields $\sim$200 million galaxy shapes --- sufficient for
  tomographic cosmic shear at $3$--$5\sigma$ DFD discrimination.
  \textbf{Impact: CRITICAL. The clock starts now: 7 months to prepare
  the analysis framework.}

\item \textbf{NEW --- Exact DFD Predictions for All Four DR1 Probes Computed
  (Section~\ref{sec:predictions}):}
  For each Euclid DR1 data product, the DFD prediction is now specified
  to formulaic precision:
  \begin{itemize}[nosep]
  \item \textbf{Cosmic shear}: $C_\ell^{\gamma\gamma}(i,j)$ suppressed by
    $[\Geff(k_\ell)/G_N]^2$ at $\ell \lesssim 100$, yielding
    $S_8^{\rm low-\ell} = 0.745 \pm 0.015$ vs
    $S_8^{\rm high-\ell} = 0.82 \pm 0.01$.
  \item \textbf{Galaxy-galaxy lensing}: Excess surface density (ESD)
    profiles show DFD $\mu$-suppression at projected separations
    $R > 10$\,Mpc, with $\Delta\Sigma_{\rm DFD}/\Delta\Sigma_\Lambda
    \approx \Geff(k \sim 1/R)/G_N$.
  \item \textbf{Galaxy clustering}: $P_{gg}(k)$ deficit of $12 \pm 4\%$
    at $k < 0.05$\,Mpc$^{-1}$ relative to \LCDM.
  \item \textbf{Photometric redshifts}: No first-order DFD effect on
    photo-$z$ estimation (distance-redshift relation unaffected at
    perturbative level). However, the psi-screen distance bias
    (Tier B) could introduce a $\sim$0.1\% systematic shift at
    $z > 1.5$ --- below DR1 detection threshold.
  \end{itemize}
  \textbf{Impact: CRITICAL. All formulae are now derivable from a single
  parameter $\kmu = 0.01$\,Mpc$^{-1}$.}

\item \textbf{NEW --- Tomographic $S_8$ Gradient Is the Killer Test:
  10 Bins, Predicted Monotonic Rise (Section~\ref{sec:tomographic}):}
  Euclid's baseline analysis uses 10 equi-populated tomographic bins with
  edges at $\{0.001, 0.418, 0.560, 0.678, 0.789, 0.900, 1.019, 1.155,
  1.324, 1.576, 2.50\}$. DFD predicts a \emph{scale-dependent} $S_8$:
  when measured from low-$\ell$ modes ($\ell < 100$, probing $k < \kmu$),
  $S_8$ is suppressed; when measured from high-$\ell$ modes
  ($\ell > 1000$), $S_8$ converges to the Planck value. The gradient
  $\Delta S_8 / \Delta \ln \ell \approx 0.03$--$0.04$ across the
  accessible multipole range is a unique DFD fingerprint that no
  $w_0 w_a$CDM, $f(R)$, or nDGP model can mimic. With
  $\sim$200 million shapes over 1900\,deg$^2$, DR1 statistical errors
  on $S_8$ per tomographic bin are $\sim$0.01--0.02, making the
  $\Delta S_8 \sim 0.07$--$0.10$ gradient a $3$--$5\sigma$ detection.
  \textbf{Impact: CRITICAL. This is the single most powerful DFD test
  in 2026.}

\item \textbf{NEW --- Weekend-Viable Minimum Analysis Pipeline Specified
  (Section~\ref{sec:weekend}):}
  A researcher can test DFD against Euclid DR1 in $\sim$2 days using
  only public data and standard tools:
  (1) Download the DR1 shear catalog and photo-$z$ catalog from
  the ESA Euclid archive;
  (2) Split into 3 broad tomographic bins ($z < 0.7$, $0.7 < z < 1.2$,
  $z > 1.2$);
  (3) Measure $C_\ell^{\gamma\gamma}$ in 8 log-spaced $\ell$-bins
  from $\ell = 50$ to $\ell = 3000$ using \texttt{NaMaster} or
  \texttt{PolSpice};
  (4) Fit $S_8$ separately for $\ell < 300$ and $\ell > 300$;
  (5) Compare $\Delta S_8$ to DFD prediction ($0.07 \pm 0.02$) vs
  \LCDM\ prediction ($0.00 \pm 0.02$).
  Total compute: $\sim$4 hours on a laptop for the power spectrum;
  $\sim$1 hour for the $S_8$ fits. The main bottleneck is downloading
  the catalog ($\sim$100\,GB).
  \textbf{Impact: HIGH. Removes all barriers to independent DFD testing.}

\item \textbf{NEW --- Competitive Landscape Mapped: Euclid MG Teams Will
  Miss the DFD Signal (Section~\ref{sec:competition}):}
  The Euclid Consortium's official modified-gravity pipeline
  (arXiv:2512.09748, arXiv:2506.03008) parameterises deviations using
  $(\mu_{\rm mg}, \Sigma_{\rm mg})$ --- phenomenological parameters that
  are \emph{scale-independent} in their baseline analysis. DFD's
  $\Geff(k)$ is intrinsically \emph{scale-dependent}. The official
  pipeline will constrain $\mu_{\rm mg}$ averaged over all scales,
  washing out the DFD signal. Specific models tested (Hu-Sawicki $f(R)$,
  nDGP, Jordan-Brans-Dicke, k-mouflage) all have different $k$-dependence
  from DFD's Yukawa form. \textbf{No team is currently testing a Yukawa
  $\Geff(k) = G_N k^2/(k^2 + \kmu^2)$ profile.} This means:
  (a) a DFD-specific analysis would be the first of its kind;
  (b) the Euclid Consortium could confirm \LCDM\ in their pipeline
  while the DFD signal hides in the scale-dependent residuals.
  \textbf{Impact: HIGH. There is a strategic window for a
  DFD-targeted Euclid DR1 paper.}

\end{enumerate}
\end{findbox}


%=============================================================================
\section{Euclid DR1 Data Products: Complete Inventory}
\label{sec:inventory}
%=============================================================================

\subsection{Release Specifications}

\begin{center}
\begin{tabular}{lll}
\toprule
\textbf{Parameter} & \textbf{DR1 Value} & \textbf{Full Survey} \\
\midrule
Release date & 21 October 2026 & $\sim$2030+ \\
Sky coverage & $\sim$1900\,deg$^2$ & 14,500\,deg$^2$ \\
VIS imaging depth (10$\sigma$ extended) & $m_{\rm AB} \sim 24.5$ & $m_{\rm AB} \sim 24.5$ \\
Galaxy shapes (WL) & $\sim$200$\times 10^6$ & $\sim$1.5$\times 10^9$ \\
Effective $n_{\rm eff}$ & $\sim$30\,arcmin$^{-2}$ & $\sim$30\,arcmin$^{-2}$ \\
Shape noise $\sigma_e$ & $\sim$0.26 & $\sim$0.26 \\
NISP photo-$z$ range & $0.2 < z < 2.5$ & $0.2 < z < 2.5$ \\
NISP spectroscopic range & $0.9 < z < 1.8$ (H$\alpha$) & same \\
Spectroscopic density & $\sim$1700\,deg$^{-2}$ & $\sim$1700\,deg$^{-2}$ \\
Photo-$z$ accuracy & $\sigma_z < 0.05(1+z)$ & $\sigma_z < 0.05(1+z)$ \\
NIR bands & $Y$, $J$, $H$ & $Y$, $J$, $H$ \\
\bottomrule
\end{tabular}
\end{center}

\subsection{Four Primary Data Products for DFD Testing}

\begin{enumerate}[label=\textbf{DP\arabic*}:]
\item \textbf{VIS Shear Catalog}: Galaxy ellipticities $e_1, e_2$ with
  PSF-corrected shapes from the VIS instrument (550--900\,nm broadband).
  This is the foundation for cosmic shear and galaxy-galaxy lensing.

\item \textbf{Photometric Redshift Catalog}: Point estimates and PDFs
  $p(z)$ for each galaxy, derived from VIS + NISP ($Y$, $J$, $H$) +
  ground-based optical photometry. Required for tomographic binning.

\item \textbf{NISP Spectroscopic Catalog}: Slitless grism spectroscopy
  providing spectroscopic redshifts for H$\alpha$ emitters at
  $0.9 < z < 1.8$. Provides the lens sample for galaxy clustering
  $P_{gg}(k)$ and cross-correlation with shear.

\item \textbf{Merged Multi-Wavelength Catalog}: Positions, magnitudes,
  morphological parameters (S\'{e}rsic fits), photometric
  classifications. Enables lens-sample selection for galaxy-galaxy
  lensing.
\end{enumerate}


%=============================================================================
\section{DFD Predictions for Each DR1 Data Product}
\label{sec:predictions}
%=============================================================================

\subsection{Cosmic Shear Power Spectrum $C_\ell^{\gamma\gamma}$}

The DFD modification enters through the growth function. In the Limber
approximation, the tomographic cosmic shear power spectrum between bins
$i$ and $j$ is:
\begin{equation}
C_\ell^{\gamma\gamma}(i,j) = \int_0^{\chi_H}
  \frac{q_i(\chi)\,q_j(\chi)}{\chi^2}\;
  P_\delta\!\left(k = \frac{\ell+1/2}{\chi},\,z(\chi)\right)\,d\chi\,,
\label{eq:Cl_shear}
\end{equation}
where $q_i(\chi)$ is the lensing kernel for bin $i$ and $P_\delta(k,z)$
is the matter power spectrum.

In DFD, the linear growth rate is modified:
\begin{equation}
\frac{d^2 D}{da^2} + \frac{3}{2a}\left(1 - \frac{w \Omega_{\rm DE}}{1+\Omega_{\rm DE}}\right)\frac{dD}{da}
  = \frac{3}{2a^2}\,\Omega_m(a)\,\frac{\Geff(k)}{G_N}\,D(a)\,,
\end{equation}
with the DFD modification:
\begin{equation}
\boxed{\frac{\Geff(k)}{G_N} = \frac{k^2}{k^2 + \kmu^2}\,,
\qquad \kmu = 0.01\;\text{Mpc}^{-1}\,.}
\label{eq:Geff}
\end{equation}

\textbf{Consequences for $C_\ell^{\gamma\gamma}$}:
\begin{itemize}[nosep]
\item At $\ell \sim 50$ (probing $k \sim 0.003$\,Mpc$^{-1}$ at median
  redshift): $\Geff/G_N \approx 0.08$. Growth is strongly suppressed.
  $C_\ell^{\gamma\gamma}$ reduced by $\sim$40--50\% relative to \LCDM.
\item At $\ell \sim 200$ ($k \sim 0.015$\,Mpc$^{-1}$):
  $\Geff/G_N \approx 0.69$. $C_\ell^{\gamma\gamma}$ reduced by
  $\sim$15--20\%.
\item At $\ell \sim 1000$ ($k \sim 0.07$\,Mpc$^{-1}$):
  $\Geff/G_N \approx 0.98$. $C_\ell^{\gamma\gamma}$ indistinguishable
  from \LCDM.
\item At $\ell > 2000$: fully \LCDM. Non-linear corrections dominate and
  DFD $\Geff$ effect vanishes.
\end{itemize}

\textbf{Prediction (E1a/E1b combined)}: Fitting $S_8$ from $\ell < 100$
yields $S_8 = 0.745 \pm 0.015$. Fitting from $\ell > 1000$ yields
$S_8 = 0.82 \pm 0.01$. The \LCDM\ prediction is $S_8 = 0.832 \pm 0.013$
at all $\ell$.


\subsection{Galaxy-Galaxy Lensing: Excess Surface Density $\Delta\Sigma(R)$}

Galaxy-galaxy lensing measures the tangential shear around lens galaxies,
yielding the excess surface density:
\begin{equation}
\Delta\Sigma(R) = \bar{\Sigma}(< R) - \Sigma(R)
  = \rho_{\rm crit}\,\Omega_m\int_0^\infty
  \left[\frac{2}{R^2}\int_0^R R'\,\xi_{gm}(\sqrt{R'^2+\Pi^2})\,dR'
  - \xi_{gm}(\sqrt{R^2+\Pi^2})\right] d\Pi\,.
\end{equation}

In DFD, the galaxy-matter cross-power spectrum is:
\begin{equation}
P_{gm}(k) = b(k)\,\frac{\Geff(k)}{G_N}\,P_{\delta}^{\Lambda}(k)\,,
\end{equation}
where $b(k)$ is the galaxy bias (assumed scale-independent on linear
scales).

\textbf{DFD $\mu$-shape in ESD profiles}:
\begin{itemize}[nosep]
\item At projected separation $R < 1$\,Mpc ($k \gg \kmu$):
  $\Delta\Sigma_{\rm DFD} \approx \Delta\Sigma_\Lambda$.
  Normal NFW-like profile.
\item At $R = 5$--$20$\,Mpc ($k \sim \kmu$):
  $\Delta\Sigma_{\rm DFD}/\Delta\Sigma_\Lambda \approx
  \Geff(k\sim 1/R)/G_N$. Deficit of 10--30\%.
\item At $R > 50$\,Mpc ($k \ll \kmu$): $\Delta\Sigma \to 0$ faster
  than \LCDM. This is the ``missing large-scale lensing'' signature.
\end{itemize}

\textbf{DR1 feasibility}: With $\sim$3 million spectroscopic lens
galaxies ($0.9 < z < 1.8$) and 200 million background shapes, the
ESD can be measured to $R \sim 30$\,Mpc with S/N $\sim$5--10 per
radial bin. The DFD deficit at $R = 10$--$30$\,Mpc is a $2$--$3\sigma$
effect in DR1 alone, strengthening with DR2.


\subsection{Galaxy Clustering $P_{gg}(k)$}

The NISP spectroscopic sample provides $P_{gg}(k)$ for H$\alpha$
emitters at $0.9 < z < 1.8$:
\begin{equation}
P_{gg}(k,z) = b^2(z)\,\left[\frac{\Geff(k)}{G_N}\right]^2
  \frac{D^2_{\rm DFD}(k,z)}{D^2_\Lambda(z)}\,P_{\rm lin}^\Lambda(k)\,,
\end{equation}
where $D_{\rm DFD}(k,z)$ is the scale-dependent DFD growth function.

\textbf{Prediction (E2)}: $P_{gg}(k < 0.05\;\text{Mpc}^{-1})$ is
suppressed by $12 \pm 4\%$ relative to the \LCDM\ prediction with the
same bias.

\textbf{Key diagnostic}: The ratio $P_{gg}(k=0.01)/P_{gg}(k=0.1)$
is lower in DFD than in \LCDM\ by a factor
$[\Geff(0.01)/\Geff(0.1)]^2 = [0.01^2/(0.01^2+0.01^2)]^2/[0.1^2/(0.1^2+0.01^2)]^2
= 0.25/0.99 \approx 0.25$. However, this applies only to the
growth-function contribution; the primordial $P(k)$ tilt partially
compensates. Net observable deficit: $\sim$12\%.

\textbf{Caveat}: Galaxy clustering alone is degenerate with scale-dependent
bias. The cross-correlation ratio $E_G(k)$ (Prediction E4) breaks this
degeneracy.


\subsection{Photometric Redshift Distribution}

\textbf{No first-order DFD effect on photo-$z$ estimation.}

The DFD $\Geff(k)$ modification affects the growth of perturbations,
not the background expansion history $H(z)$ or the luminosity-distance
relation $d_L(z)$. Photo-$z$ estimation uses galaxy SEDs and apparent
magnitudes, which depend on $d_L(z)$ and intrinsic luminosity --- neither
is modified by DFD at the perturbative level.

\textbf{Second-order effect (Tier B, psi-screen)}: If the DFD
psi-screen mechanism introduces a distance bias
$\delta d_L/d_L \sim 10^{-4}$--$10^{-3}$ at $z > 1.5$, this could
shift mean photo-$z$ by $\delta z \sim 0.001$--$0.01$. This is below
the Euclid photo-$z$ accuracy requirement of $0.002(1+z)$ and
undetectable in DR1. It becomes relevant only with full-survey
statistics and external spectroscopic calibration.


%=============================================================================
\section{Tomographic $S_8$: The Killer Test}
\label{sec:tomographic}
%=============================================================================

\subsection{Euclid Tomographic Bin Structure}

The Euclid baseline cosmic shear analysis uses \textbf{10 equi-populated
tomographic bins} with bin edges:
\begin{equation}
z_{\rm edges} = \{0.001,\; 0.418,\; 0.560,\; 0.678,\; 0.789,\; 0.900,\;
  1.019,\; 1.155,\; 1.324,\; 1.576,\; 2.50\}\,.
\end{equation}

With 1900\,deg$^2$ and 30\,arcmin$^{-2}$, each bin contains
$\sim$20 million galaxies. The auto- and cross-power spectra form a
$10 \times 10$ matrix of 55 independent spectra $C_\ell^{ij}$.

\subsection{DFD Predicted Gradient}

In \LCDM, $S_8$ is scale-independent: fitting from any $\ell$-range
or any combination of tomographic bins yields the same $S_8$ (up to
noise and systematics).

In DFD, $S_8$ is \emph{scale-dependent} because $\Geff(k)$ modifies
the growth function differently at each $k$. The observable $S_8$
depends on which $\ell$-range dominates the fit:

\begin{center}
\begin{tabular}{lccc}
\toprule
\textbf{$\ell$-range} & \textbf{Effective $k$ (Mpc$^{-1}$)}
  & \textbf{$S_8^{\rm DFD}$} & \textbf{$S_8^{\Lambda{\rm CDM}}$} \\
\midrule
$50 < \ell < 100$   & $0.003$--$0.007$ & $0.73 \pm 0.02$ & $0.832 \pm 0.013$ \\
$100 < \ell < 300$  & $0.007$--$0.02$  & $0.76 \pm 0.015$ & $0.832 \pm 0.013$ \\
$300 < \ell < 1000$ & $0.02$--$0.07$   & $0.80 \pm 0.01$ & $0.832 \pm 0.013$ \\
$1000 < \ell < 3000$ & $0.07$--$0.2$   & $0.82 \pm 0.01$ & $0.832 \pm 0.013$ \\
\bottomrule
\end{tabular}
\end{center}

\textbf{Key diagnostic}: Plot $S_8(\ell_{\rm max})$ --- the value of
$S_8$ obtained by fitting only $C_\ell^{ij}$ up to $\ell_{\rm max}$.
DFD predicts a \emph{monotonically rising} curve from
$S_8 \approx 0.73$ at $\ell_{\rm max} = 100$ to $S_8 \approx 0.82$ at
$\ell_{\rm max} = 3000$. \LCDM\ predicts a flat line at $0.832$.

\textbf{Discrimination}: The gradient $\Delta S_8 \equiv
S_8(\ell > 1000) - S_8(\ell < 100) = 0.07$--$0.10$ in DFD vs $0.00$
in \LCDM. With DR1 errors of $\sim$0.02 per $\ell$-band, this is a
$3$--$5\sigma$ detection.

\textbf{Why this is unique to DFD}: Alternative explanations for the
$S_8$ tension ($w_0 w_a$CDM, massive neutrinos, baryonic feedback)
suppress power at high $k$ (small scales), which would produce the
\emph{opposite} gradient: $S_8$ decreasing with $\ell$. Only DFD's
large-scale suppression via $\Geff(k \ll \kmu) \to 0$ produces the
low-$\ell$ deficit / high-$\ell$ agreement pattern.


%=============================================================================
\section{Analysis Pipeline: Shear Catalog to DFD Discrimination}
\label{sec:pipeline}
%=============================================================================

\begin{pipelinebox}[Step-by-Step Analysis Pipeline]

\textbf{Inputs}: Euclid DR1 shear catalog (galaxy positions, $e_1$, $e_2$,
weight, photo-$z$ PDF).

\medskip
\textbf{Step 1 --- Catalog preparation} ($\sim$2 hours):
\begin{itemize}[nosep]
\item Download from ESA Euclid Science Archive
  (\texttt{https://easidr.esac.esa.int/}) or NASA/IPAC IRSA archive.
\item Apply recommended quality cuts (S/N, size, flag selections per
  Euclid Consortium recommendations).
\item Assign galaxies to tomographic bins using photo-$z$ point estimates
  or stacked $p(z)$.
\end{itemize}

\textbf{Step 2 --- Mask and survey geometry} ($\sim$1 hour):
\begin{itemize}[nosep]
\item Download the DR1 survey mask (binary HEALPix or polygon format).
\item Compute the effective area and identify masked regions (bright
  stars, satellite trails, CCD gaps).
\item Compute mixing matrices for pseudo-$C_\ell$ correction.
\end{itemize}

\textbf{Step 3 --- Power spectrum estimation} ($\sim$4 hours):
\begin{itemize}[nosep]
\item Use \texttt{NaMaster} (Alonso et al.\ 2019) or \texttt{PolSpice}
  for pseudo-$C_\ell$ estimation with mode-coupling correction.
\item Measure all 55 auto- and cross-spectra $C_\ell^{ij}$ in
  $\sim$15--20 logarithmic $\ell$-bins from $\ell = 30$ to $\ell = 5000$.
\item Subtract noise bias: $N_\ell^{ii} = \sigma_e^2/\bar{n}_i$ for
  auto-spectra.
\item Estimate covariance via Gaussian analytic formula (sufficient for
  DR1) or jackknife over sub-patches.
\end{itemize}

\textbf{Step 4 --- DFD model fitting} ($\sim$2 hours):
\begin{itemize}[nosep]
\item Implement the DFD-modified $C_\ell^{ij}$ theory using a Boltzmann
  code (\texttt{CLASS} or \texttt{CAMB}) with $\Geff(k)$ modification
  (Eq.~\ref{eq:Geff}) injected into the growth equation.
\item Alternatively, use the approximation:
  $C_\ell^{ij,\rm DFD} \approx [\Geff(k_\ell^{\rm eff})/G_N]^2 \times
  C_\ell^{ij,\Lambda}$, where $k_\ell^{\rm eff} \approx
  (\ell + 1/2)/\chi(z_{\rm eff})$.
\item Run MCMC (e.g.\ \texttt{emcee}, \texttt{cobaya}) sampling
  $\{S_8,\,\Omega_m,\,\kmu\}$ with \LCDM\ parameters fixed to Planck
  2018 values.
\item Compare: nested model selection with $\kmu = 0$ (\LCDM) vs
  $\kmu$ free (DFD). Compute Bayes factor or $\Delta\chi^2$.
\end{itemize}

\textbf{Step 5 --- Scale-dependent $S_8$ diagnostic} ($\sim$1 hour):
\begin{itemize}[nosep]
\item Fit $S_8$ separately in 4 $\ell$-bands: [50, 100], [100, 300],
  [300, 1000], [1000, 3000].
\item Plot $S_8(\ell\text{-band})$ with error bars.
\item DFD signature: monotonic rise. \LCDM: flat.
\item Quantify with a linear fit: $S_8 = \alpha + \beta \ln(\ell)$.
  DFD predicts $\beta = 0.010 \pm 0.003$; \LCDM\ predicts $\beta = 0$.
\end{itemize}

\textbf{Step 6 --- Cross-checks and systematics} ($\sim$4 hours):
\begin{itemize}[nosep]
\item B-mode check: $C_\ell^{BB}$ should be consistent with zero
  (systematic contamination test).
\item PSF residual test: correlate residual PSF ellipticity with shear.
\item Photo-$z$ calibration: compare $n(z)$ from stacked $p(z)$ with
  spectroscopic sub-sample cross-correlation.
\item Baryonic feedback test: vary the \texttt{HMCode} baryon parameter
  $A_{\rm bary}$ and check that DFD signal persists.
\end{itemize}

\textbf{Total wall-clock time}: $\sim$14 hours for an experienced
analyst. $\sim$2--3 days including setup and debugging.

\end{pipelinebox}


%=============================================================================
\section{Weekend-Viable Minimum Analysis}
\label{sec:weekend}
%=============================================================================

\begin{granitebox}[CAN SOMEONE TEST DFD IN A WEEKEND? YES.]

\textbf{Prerequisites}:
\begin{itemize}[nosep]
\item Python 3.10+ with \texttt{numpy}, \texttt{scipy}, \texttt{healpy},
  \texttt{pymaster} (NaMaster), \texttt{emcee}.
\item Laptop with 16\,GB RAM (sufficient for 1900\,deg$^2$ catalog at
  HEALPix $N_{\rm side} = 1024$).
\item Internet connection for catalog download ($\sim$50--100\,GB;
  pre-download Friday evening).
\end{itemize}

\textbf{Saturday morning (4 hours)}:
\begin{enumerate}[nosep]
\item Load shear catalog, apply quality cuts.
\item Split into 3 broad tomographic bins:
  $z < 0.7$ ($\sim$60M galaxies), $0.7 < z < 1.2$ ($\sim$80M),
  $z > 1.2$ ($\sim$60M).
\item Pixelise shear field onto HEALPix maps ($N_{\rm side} = 1024$,
  $\sim$3.4\,arcmin pixels).
\item Measure 6 $C_\ell^{ij}$ spectra in 8 log-spaced $\ell$-bins
  ($\ell \in [50, 100, 200, 400, 800, 1500, 3000, 5000]$).
\end{enumerate}

\textbf{Saturday afternoon (3 hours)}:
\begin{enumerate}[nosep, start=5]
\item Compute Gaussian covariance matrix analytically.
\item Fit $S_8$ from $\ell < 300$ modes only (DFD-sensitive range).
\item Fit $S_8$ from $\ell > 300$ modes only (DFD-insensitive range).
\item Compare: $\Delta S_8 = S_8^{\rm high-\ell} - S_8^{\rm low-\ell}$.
\end{enumerate}

\textbf{Sunday (4 hours)}:
\begin{enumerate}[nosep, start=9]
\item Run MCMC with $\kmu$ as free parameter: flat prior
  $\kmu \in [0, 0.1]$\,Mpc$^{-1}$.
\item Compute posterior on $\kmu$ and Bayes factor vs \LCDM.
\item B-mode null test.
\item Write up result.
\end{enumerate}

\textbf{Expected outcome}:
\begin{itemize}[nosep]
\item If DFD correct: $\Delta S_8 = 0.07 \pm 0.02$, $\kmu$ posterior
  peaked at $\sim$0.01\,Mpc$^{-1}$, Bayes factor $> 10$ favouring DFD.
\item If \LCDM\ correct: $\Delta S_8 = 0.00 \pm 0.02$, $\kmu$ posterior
  consistent with zero, Bayes factor $< 1$.
\end{itemize}

\end{granitebox}


%=============================================================================
\section{Competitive Landscape}
\label{sec:competition}
%=============================================================================

\subsection{Who Will Analyze Euclid DR1 for Modified Gravity?}

\begin{enumerate}[label=\textbf{\arabic*.}]

\item \textbf{Euclid Consortium Internal (Theory Working Group, SWG-TH)}:
  $\sim$200 scientists dedicated to beyond-\LCDM\ analysis.
  \begin{itemize}[nosep]
  \item \textbf{Methods}: $(\mu_{\rm mg}, \Sigma_{\rm mg})$
    parameterisation (scale-independent baseline), effective field
    theory (EFT) of dark energy, specific models ($f(R)$, nDGP,
    k-mouflage, JBD).
  \item \textbf{DFD blind spot}: The baseline $(\mu_{\rm mg},
    \Sigma_{\rm mg})$ analysis averages over scales, washing out the
    $k$-dependent DFD signal. Their specific-model tests (Hu-Sawicki,
    nDGP) have different $k$-dependence (chameleon screening in $f(R)$,
    Vainshtein screening in nDGP) and will not detect a Yukawa form.
  \end{itemize}

\item \textbf{External Survey Teams (DES, KiDS, HSC)}:
  Will cross-analyse Euclid DR1 with their own shear catalogs for
  consistency checks.
  \begin{itemize}[nosep]
  \item Focus is on $S_8$ tension resolution, not specific MG models.
  \item They will measure a single $S_8$ from the full $\ell$-range,
    potentially averaging out the DFD signal.
  \end{itemize}

\item \textbf{DESI Cross-Correlation Teams}:
  Will combine Euclid lensing with DESI spectroscopy for $E_G(k)$ and
  $f\sigma_8$ measurements.
  \begin{itemize}[nosep]
  \item Most likely to detect the DFD signal accidentally, since
    $E_G(k)$ is explicitly scale-dependent.
  \item However, typical $E_G$ analyses use broad $k$-bins, which may
    still miss the narrow Yukawa transition.
  \end{itemize}

\item \textbf{Model-Independent Reconstruction Groups}:
  Principal component analyses of $\mu(k,z)$ and $\Sigma(k,z)$.
  \begin{itemize}[nosep]
  \item Most sensitive to DFD-type signals, but typically lack the
    statistical power with DR1 alone to constrain more than 2--3 PCs.
  \item Paper by arXiv:2503.20951 (March 2025) explores model-agnostic
    $3\times 2$pt PCA --- closest existing framework to detecting DFD.
  \end{itemize}

\end{enumerate}

\subsection{How DFD Analysis Differs}

\begin{center}
\begin{tabular}{lcc}
\toprule
\textbf{Feature} & \textbf{Standard MG} & \textbf{DFD} \\
\midrule
$\Geff$ parameterisation & $\mu(a)$ (scale-free) & $k^2/(k^2+\kmu^2)$ \\
Number of free parameters & 2 ($\mu_0$, $\eta_0$) & 1 ($\kmu$) \\
Screening mechanism & Chameleon/Vainshtein & None (Yukawa cutoff) \\
Transition scale & Absent or environment-dep. & $\kmu^{-1} = 100$\,Mpc \\
$S_8$ prediction & Scale-independent shift & Scale-dependent gradient \\
Key observable & Amplitude change & $\ell$-dependent $S_8$ \\
\bottomrule
\end{tabular}
\end{center}

\textbf{Strategic implication}: A DFD-targeted paper analyzing Euclid
DR1 with the specific Yukawa ansatz would be the \emph{first} such
analysis. The one-parameter nature ($\kmu$ only) makes it trivially
implementable and publishable. No other team is currently pursuing this.


%=============================================================================
\section{Euclid Falsification Contingency}
\label{sec:euclid_contingency}
%=============================================================================

\begin{contigbox}[Contingency: What If Euclid DR1 Shows No $S_8$ Gradient?]

Building on the three-level contingency from W4i16 (developed for SO),
the Euclid-specific contingency:

\textbf{Level 1 --- $\Delta S_8 < 0.04$ ($< 2\sigma$ tension with DFD)}:
\begin{itemize}[nosep]
\item Consistent with noise. DR1 covers only 1900\,deg$^2$ (13\% of
  full survey). Low-$\ell$ $S_8$ has large cosmic variance.
\item \textbf{Action}: Wait for DR2 ($\sim$7000\,deg$^2$, expected 2028),
  which reduces $\ell < 100$ errors by $\sim$2$\times$.
\item \textbf{Status}: DFD fully viable. No adjustment needed.
\end{itemize}

\textbf{Level 2 --- $\Delta S_8 = 0.00 \pm 0.02$ ($3$--$4\sigma$ against DFD)}:
\begin{itemize}[nosep]
\item DFD $\kmu = 0.01$\,Mpc$^{-1}$ is disfavoured but not excluded.
\item \textbf{Action}: Shift to $\kmu \leq 0.005$\,Mpc$^{-1}$ (crossover
  at $> 200$\,Mpc). This preserves the framework but pushes the observable
  effect to $\ell < 30$, where DR1 has no statistical power but DR2+
  and SO will probe.
\item \textbf{Status}: DFD survives in squeezed form.
\end{itemize}

\textbf{Level 3 --- $\Delta S_8 = 0.00 \pm 0.01$ with DR2+ confirmation
  ($> 5\sigma$ against DFD)}:
\begin{itemize}[nosep]
\item The Yukawa $\Geff(k)$ form is falsified at $k > 0.003$\,Mpc$^{-1}$.
\item \textbf{Action}: Same as W4i16 Level 3 --- retreat to galactic-scale
  DFD. The refractive framework survives; the cosmological $\Geff$
  sector is abandoned.
\item \textbf{Status}: Major reformulation.
\end{itemize}

\textbf{Key point}: DR1 alone cannot reach Level 3 due to limited area.
The worst DR1 can do is Level 2, which is survivable. The decisive test
requires DR2 (2028) or full survey.

\end{contigbox}


%=============================================================================
\section{Timeline and Action Items}
\label{sec:timeline}
%=============================================================================

\begin{center}
\begin{tabular}{llp{7cm}}
\toprule
\textbf{Date} & \textbf{Milestone} & \textbf{DFD Action} \\
\midrule
Now (Mar 2026) & Q1 data public & Practice pipeline on 63\,deg$^2$ Q1 data.
  Test shear measurement, photo-$z$ cuts, power spectrum estimation. \\
May 2026 & --- & Implement $\Geff(k)$ in \texttt{CLASS} or build
  lookup-table approximation. Validate against existing $S_8$ measurements
  (DES Y6, KiDS-Legacy). \\
Aug 2026 & --- & Pre-compute DFD $C_\ell^{ij}$ theory vectors for
  $\kmu \in [0.001, 0.1]$ grid. Prepare MCMC sampler. Draft paper
  skeleton. \\
21 Oct 2026 & \textbf{Euclid DR1} & Execute pipeline. Target 2-week
  turnaround for first result. \\
Nov 2026 & First result & Internal DFD assessment: $\Delta S_8$ consistent
  with DFD or \LCDM? \\
Dec 2026 & Paper submission & ``Testing Density Field Dynamics against
  Euclid DR1: A scale-dependent $S_8$ analysis'' \\
Q4 2026--Q2 2027 & SO lensing & Cross-check with CMB lensing
  $C_\ell^{\kappa\kappa}$. \\
\bottomrule
\end{tabular}
\end{center}

\textbf{Immediate priority}: Use the Euclid Q1 data (63\,deg$^2$,
released March 2025, publicly available) as a testbed. While Q1
lacks the area for cosmological constraints, it allows end-to-end
pipeline validation: shear catalog handling, tomographic binning,
$C_\ell$ estimation, and systematic checks.


\end{document}
